%==============================================================================
% 四、符号说明
%==============================================================================
\section{符号说明}

为便于对照，全文主要符号列于表~\ref{tab:symbols}。未在此表中的一次性符号在首次出现处就地说明。

\begin{table}[htbp]
\centering
\small
\caption{主要符号说明}
\label{tab:symbols}
\begin{tabular*}{\textwidth}{@{\extracolsep{\fill}}p{1.5cm}p{9.3cm}p{2.6cm}@{}}
\toprule
\textbf{符号} & \textbf{含义} & \textbf{单位} \\
\midrule
$n,\ p,\ m$ & 样本数、化学特征维数、感官维数，其中 $n=250$、$p=231$、$m=50$ & -- \\
$c_{ij}$ & 第 $i$ 个啤酒第 $j$ 个化合物的浓度，已还原量纲 & mg/L 等 \\
$y_{is}$ & 第 $i$ 个啤酒第 $s$ 个感官维度的强度值 & 无量纲 \\
$X,\ Y$ & 化学特征矩阵与感官标签矩阵 & -- \\
$\phi_s$ & 第 $s$ 个感官维度的左尾占比，即与该列最小值相距 ≤0.15 的样本占比 & -- \\
$\pi_j$ & 化合物或化学族 $j$ 的置换重要性 & -- \\
$\pi_{\mathrm{thr}}$ & 置换重要性的入选频率阈值（非 FDR 控制值） & -- \\
$B,\ k$ & 自助重采样次数与 Top-$k$ 入选集合，其中 $B=200$、$k=3$ & -- \\
$O_{ij}$ & 化合物 $j$ 对感官维度 $i$ 的气味活性值 OAV & -- \\
$J(A,B)$ & 两个候选集合的 Jaccard 系数 & -- \\
$x_j,\ x^{0}_j$ & 问题三中化合物 $j$ 的优化值与原始值（$\log_{10}$ 域） & -- \\
$z_j$ & 化合物 $j$ 是否被改动的 0--1 决策变量 & -- \\
$\hat y_s(x)$ & 线性代理对第 $s$ 个感官维度的预测 & 无量纲 \\
$t_s$ & 目标风格的参考感官值（风格内中位数） & 无量纲 \\
$d_k,\ \tau_k$ & 缺陷物 $k$ 的浓度与其香气阈值 & $\mu$g/L 或 mg/L \\
$f_1,f_2,f_3$ & 感官偏差、缺陷风险、改动代价三个目标 & 无量纲 \\
$\varepsilon_2,\ \varepsilon_3$ & $\varepsilon$-约束法中 $f_2,f_3$ 的上界 & 无量纲 \\
\bottomrule
\end{tabular*}
\end{table}
